% =============================================================================
%  Grundlagen der Signalverarbeitung Latex Skript
%  von Kaan Büyüksahin
%  letzte Aktualisieurng 05. September 2025
% =============================================================================

\documentclass[12pt]{article}
\usepackage[rgb, table]{xcolor}
\usepackage{pgfplots}
\usepackage{graphicx}
\usepackage{svg}
\usepackage{amsmath}
\usepackage{afterpage}
\usepackage{fancyhdr}
\usepackage{multicol}
\usepackage{tikz}
\usepackage{transparent}
\usepackage{amssymb}
\usetikzlibrary{trees, shadings}
\usepackage{caption}
\usepackage{tabularx}
\usepackage{circuitikz}
\usepackage{geometry}
\geometry{
a4paper,
left = 1cm,
right = 1cm,
top = 1 cm,
bottom = 1 cm
}
\pgfplotsset{compat=1.15}
\usepackage[ngerman]{babel}
\usepackage{datetime}
\renewcommand{\dateseparator}{/}
\usepackage{times}
\usepackage[utf8]{inputenc}
\usepackage[T1]{fontenc}
\usepackage[
   hidelinks,
   pdfauthor={Kaan Büyükşahin},
   pdftitle={Grundlagen der Signalverarbeitung - Kurzreferenz}, % chktex 8
   pdfsubject={Zusammenfassung der wichtigsten Inhalte für Grundlagen der Signalverarbeitung},
   pdfcreator={pdflatex},
   ]{hyperref}

%%% Farbsettings %%%
\definecolor{Aufgabe1Teil1}{rgb}{0.929, 1.0, 0.859}
\definecolor{Aufgabe1Teil2}{rgb}{0.894, 0.953, 0.894}
\definecolor{Aufgabe2}{rgb}{1.0, 0.796, 0.796}
\definecolor{Aufgabe4}{rgb}{0.910, 0.961, 0.988}

% TikZ Bibliotheken für Farbverläufe
\usetikzlibrary{shadings}

% Benutzerdefinierte Farben definieren
\definecolor{color1}{RGB}{82, 57, 128}
\definecolor{color2}{RGB}{196,173,221}
\definecolor{color3}{RGB}{27, 14, 32}
\definecolor{color4}{RGB}{200, 179, 246}
\definecolor{color5}{RGB}{209, 192, 236}

% Titel-Informationen
\title{Grundlagen der Signalverarbeitung}
\author{Kaan Büyükşahin}
\date{\today}

% Kopf- und Fußzeile konfigurieren
\pagestyle{fancy}
\fancyhf{}
\rhead{\thepage}
\lhead{Grundlagen der Signalverarbeitung}
\renewcommand{\headrulewidth}{0.4pt}

\begin{document}

\pagestyle{fancy}
\fancyhf{}
\fancyfoot[R]{\raisebox{10pt}{\thepage}} % Seitenzahl nach oben gezogen
\renewcommand{\headrulewidth}{0pt}  % Unteren Balken gelsöcht

%%% Deckblatt %%%

% Titelseite mit Farbverlauf
\begin{titlepage}
    \begin{tikzpicture}[remember picture,overlay]
        % Haupthintergrund mit sanftem diagonalem Farbverlauf
        \shade[shading=linear,left color=color1,right color=color4] 
            (current page.south west) rectangle (current page.north east);

      % Zweite Schicht für weichere Übergänge (vertikal)
        \shade[shading=linear,top color=color2!60,bottom color=color3,opacity=0.3] 
            (current page.south west) rectangle (current page.north east);
         
    \end{tikzpicture}

    % ZENTRIERUNG: Alles wird horizontal zentriert
    \centering
    
    % ABSTAND VON OBEN: 3.5cm Platz vom Seitenrand
    \vspace*{3.5cm}
    
    % HAUPTTITEL - eleganter und moderner
    {\fontsize{40}{48}\selectfont\transparent{0.85}\color{white!}\bfseries 
    Grundlagen der\\[0.4cm] 
    \transparent{0.85}\textcolor{white!}{Signalverarbeitung}}
    
    % ABSTAND ZWISCHEN TITEL UND UNTERTITEL: 2cm
    \vspace{2cm}
    
    % UNTERTITEL - stilvoller 
    {\fontsize{22}{26}\selectfont\transparent{0.85}\color{white}  
    Kurzreferenz}
    
    % FLEXIBLER PLATZ: Drückt den unteren Bereich nach unten
    \vspace*{\fill}

    % INFORMATIONSBEREICH - professioneller gestaltet
    \transparent{0.85}\color{white}
    \begin{minipage}[t]{0.55\linewidth}
        \raggedright{}
        Kaan Büyükşahin\\
        Technische Universität Darmstadt\\
        Fachbereich Elektro- und Informationstechnik\\
        Fachgebiet Signalverarbeitung bei Prof.\ Dr.-Ing.\ Zoubir
    \end{minipage}
    \hfill
    \begin{minipage}[t]{0.40\linewidth}
        \raggedleft{}
        {\href{https://www.kaanbs.com}{\transparent{0.85}\textcolor{white}{kaanbs.com}}}\\
        Aktualisiert: {\normalsize\today}
    \end{minipage}
    
    % ABSTAND ZUM SEITENRAND: 2cm vom unteren Rand
    \vspace{1cm}

\end{titlepage}

%%% Mathematische Zusammenhänge ===================================================================

\section*{\underline{Mathematische Zusammenhänge}}

   \subsection*{Trigonometrische Funktionen}
      \subsubsection*{Winkeltabelle}

         \begin{table}[h!]
            \centering
            \renewcommand{\arraystretch}{2} 
            \resizebox{\textwidth}{!}{
            \scriptsize
               \begin{tabularx}{\textwidth}{|c|*{17}{>{\centering\arraybackslash}X|}}  % chktex 44
                  \hline   % chktex 44
                  \textbf{DEG} & \(0^{\circ}\) & \(30^{\circ}\) & \(45^{\circ}\) & \(60^{\circ}\) & \(90^{\circ}\) & \(120^{\circ}\) & \(135^{\circ}\) & \(150^{\circ}\) & \(180^{\circ}\) & \(210^{\circ}\) & \(225^{\circ}\) & \(240^{\circ}\) & \(270^{\circ}\) & \(300^{\circ}\) & \(315^{\circ}\) & \(330^{\circ}\) & \(360^{\circ}\) \\
                  \hline   % chktex 44
                  \textbf{rad} & \(0\) & \(\frac{\pi}{6}\) & \(\frac{\pi}{4}\) & \(\frac{\pi}{3}\) & \(\frac{\pi}{2}\) & \(\frac{2\pi}{3}\) & \(\frac{3\pi}{4}\) & \(\frac{5\pi}{6}\) & \(\pi\) & \(\frac{7\pi}{6}\) & \(\frac{5\pi}{4}\) & \(\frac{4\pi}{3}\) & \(\frac{3\pi}{2}\) & \(\frac{5\pi}{3}\) & \(\frac{7\pi}{4}\) & \(\frac{11\pi}{6}\) & \(2\pi\) \\
                  \hline   % chktex 44
                  \hline   % chktex 44
                  \(\sin{x}\) & \(0\) & \(\frac{1}{2}\) & \(\frac{\sqrt{2}}{2}\) & \(\frac{\sqrt{3}}{2}\) & \(1\) & \(\frac{\sqrt{3}}{2}\) & \(\frac{\sqrt{2}}{2}\) & \(\frac{1}{2}\) & \(0\) & \(-\frac{1}{2}\) & \(-\frac{\sqrt{2}}{2}\) & \(-\frac{\sqrt{3}}{2}\) & \(-1\) & \(-\frac{\sqrt{3}}{2}\) & \(-\frac{\sqrt{2}}{2}\) & \(-\frac{1}{2}\) & \(0\) \\
                  \hline   % chktex 44
                  \(\cos{x}\) & \(1\) & \(\frac{\sqrt{3}}{2}\) & \(\frac{\sqrt{2}}{2}\) & \(\frac{1}{2}\) & \(0\) & \(-\frac{1}{2}\) & \(-\frac{\sqrt{2}}{2}\) & \(-\frac{\sqrt{3}}{2}\) & \(-1\) & \(-\frac{\sqrt{3}}{2}\) & \(-\frac{\sqrt{2}}{2}\) & \(-\frac{1}{2}\) & \(0\) & \(\frac{1}{2}\) & \(\frac{\sqrt{2}}{2}\) & \(\frac{\sqrt{3}}{2}\) & \(1\) \\
                  \hline   % chktex 44
                  \(\tan{x}\) & \(0\) & \(\frac{1}{\sqrt{3}}\) & \(1\) & \(\sqrt{3}\) & n.\ d. & \(-\sqrt{3}\) & \(-1\) & \(-\frac{1}{\sqrt{3}}\) & \(0\) & \(\frac{1}{\sqrt{3}}\) & \(1\) & \(\sqrt{3}\) & n.\ d. & \(-\sqrt{3}\) & \(-1\) & \(-\frac{1}{\sqrt{3}}\) & \(0\) \\
                  \hline   % chktex 44
               \end{tabularx} 
            }
         \end{table}

      \smallskip

      \noindent
      \begin{minipage}[t]{0.45\textwidth}
         \subsubsection*{Pythagoras}
            \[
               \cos^2(x) + \sin^2(x) = 1
            \]

         \subsubsection*{Symmetrie und Verschiebung}
            \begin{align*}
               \cos(-x) &= \cos(x) \quad \text{(gerade Funktion)} \\
               \sin(-x) &= -\sin(x) \quad \text{(ungerade Funktion)} \\
               \sin\left(x + \frac{\pi}{2}\right) &= \cos(x) \\
               \cos\left(x - \frac{\pi}{2}\right) &= \sin(x)
            \end{align*}

         \subsubsection*{Doppelwinkelfunktionen}
            \begin{align*}
               \cos^2(x) &= \frac{1}{2}\left(1 + \cos(2x)\right) \\
               \sin^2(x) &= \frac{1}{2}\left(1 - \cos(2x)\right) \\
               \cos(2x) &= \cos^2(x) - \sin^2(x) \\ &= 1 - 2\sin^2(x) = 2\cos^2(x) - 1 \\
               \sin(2x) &= 2\sin(x)\cos(x)
            \end{align*}
      \end{minipage}
      \hfill
      \begin{minipage}[t]{0.45\textwidth}
         \subsubsection*{Additionstheoreme}
            \begin{align*}
               \cos(x \pm y) &= \cos(x)\cos(y) \mp \sin(x)\sin(y) \\
               \sin(x \pm y) &= \sin(x)\cos(y) \pm \cos(x)\sin(y) \\
               \sin(x) + \sin(y) &= 2\sin\left(\frac{x+y}{2}\right)\cos\left(\frac{x-y}{2}\right) \\
               \sin(x) - \sin(y) &= 2\cos\left(\frac{x+y}{2}\right)\sin\left(\frac{x-y}{2}\right) \\
               \sin(x)\sin(y) &= \frac{1}{2}\left[\cos(x-y) - \cos(x+y)\right] \\
               \cos(x) + \cos(y) &= 2\cos\left(\frac{x+y}{2}\right)\cos\left(\frac{x-y}{2}\right) \\
               \cos(x) - \cos(y) &= -2\sin\left(\frac{x+y}{2}\right)\sin\left(\frac{x-y}{2}\right) \\
               \cos(x)\cos(y) &= \frac{1}{2}\left[\cos(x-y) + \cos(x+y)\right] \\
               \sin(x)\cos(y) &= \frac{1}{2}\left[\sin(x-y) + \sin(x+y)\right]
            \end{align*}
      \end{minipage}

      \newpage
      
   \subsection*{Komplexe Zahlen und Integrale}
      \subsubsection*{Komplexer Zeiger mit Einheitskreis und Punkten}

         \begin{center}
            \begin{tabular}{c c} 
            \raisebox{3 cm}{ 
               \begin{minipage}[t]{0.45\textwidth}
                  \renewcommand{\arraystretch}{1.5}
                  \begin{tabular}{|c|c|c|c|c|c|}   % chktex 44
                  \hline   % chktex 44
                  \textbf{\(\varphi\)} & \textbf{\(0\)} & \textbf{\(\frac{\pi}{2}\)} & \textbf{\(\pi\)} & \textbf{\(\frac{3\pi}{2}\)} & \textbf{\(2\pi\)} \\
                  \hline   % chktex 44
                  \(x = e^{j\varphi}\) & \(1\) & \(j\) & \(-1\) & \(-j\) & \(1\) \\
                  \hline   % chktex 44
                  Punkt & \((1,0)\) & \((0,1)\) & \((-1,0)\) & \((0,-1)\) & \((1,0)\) \\
                  \hline   % chktex 44
                  \noalign{\vskip 10pt}
                  \hline   % chktex 44
                  \(-\varphi\) & \textbf{\(0\)} & \textbf{\(-\frac{\pi}{2}\)} & \textbf{\(-\pi\)} & \textbf{\(-\frac{3\pi}{2}\)} & \textbf{\(-2\pi\)} \\
                  \hline   % chktex 44
                  \(x = e^{-j\varphi}\) & \(1\) & \(-j\) & \(-1\) & \(j\) & \(1\) \\
                  \hline   % chktex 44
                  Punkt & \((1,0)\) & \((0,-1)\) & \((-1,0)\) & \((0,1)\) & \((1,0)\) \\
                  \hline   % chktex 44
               \end{tabular}
               \end{minipage}
               }
               &
               \begin{minipage}[t]{0.55\textwidth} 
                  \begin{center}
                     \begin{tikzpicture}[scale=1.8] 
                     % Der Kreis
                     \draw[thick] (1,0) circle(1.25); 

                     % Achsen
                     \draw[->] (-0.5,0) -- (2.5,0) node[anchor=west] {\text{Re}(x)};
                     \draw[->] (1,-1.5) -- (1,1.55) node[anchor=south] {\text{Im}(x)};

                     % Punkte auf dem Einheitskreis
                     \filldraw[blue] (2.25 , 0) circle(0.04);
                     \filldraw[blue] (1 , 1.25) circle(0.04); 
                     \filldraw[blue] (-0.25 , 0) circle(0.04);
                     \filldraw[blue] (1 , -1.25) circle(0.04);

                     % Winkelbeschriftungen
                     \node at (2.35, -0.2) {\(1\)};
                     \node at (0.85, 1.425) {\(1\)};
                     \node at (-0.45, -0.2) {\(-1\)};
                     \node at (0.8, -1.45) {\(-1\)};
                     \end{tikzpicture}
                  \end{center}
               \end{minipage}
            \end{tabular}
      \end{center}


      \begin{minipage}[t]{0.48\textwidth}
      \subsubsection*{Grundlagen komplexer Zahlen}
         \begin{align*}
            j^2 &= -1 \\
            -j^2 &= 1 \\
            j &= \sqrt{-1} \\
            j^{-1} &= \frac{1}{j} = -j \\
            j^3 &= j \cdot j^2 = -j \\
            j^4 &= (j^2)^2 = 1   % chktex 3
         \end{align*}    

      \subsubsection*{Darstellung von Sinus und Kosinus}
         \begin{align*}
            e^{j\varphi} &= \cos(\varphi) + j\sin(\varphi) \\
            \cos(\varphi) &= \frac{1}{2}\left(e^{j\varphi} + e^{-j\varphi}\right) \\
            \sin(\varphi) &= \frac{1}{2j}\left(e^{j\varphi} - e^{-j\varphi}\right)
         \end{align*}
      \subsection*{Betragsdefinitionen}
         \begin{align*}
            |x| &= \sqrt{x^2} \quad (x \in \mathbb{R}) \\
            |z| &= \sqrt{\mathrm{Re}(z)^2 + \mathrm{Im}(z)^2} \quad (z \in \mathbb{C}) \\ % chktex 3
            z &= x + j\,y, \\
            |z|^2 &= z \cdot z^* = \mathrm{Re}(z)^2 + \mathrm{Im}(z)^2 \\  % chktex 3
            |e^{\,j\varphi}| &= 1 \\
            |e^{-\,j\varphi}| &= 1 \\
         \end{align*}

      \end{minipage}
      \hfill
      \begin{minipage}[t]{0.48\textwidth}
         \subsection*{Integrale}

            \noindent
            \(\int \frac{1}{x^n} \, dx = \frac{x^{1-n}}{1-n} + C, \quad \text{für } n \neq 1\) \\[8pt]
            \(\int x \sqrt{X} \, dx = -\frac{1}{3} \sqrt{X^3} \quad \text{mit } X = a^2 - x^2\)\\[8pt]
            \(\int x^2 \sqrt{X} \, dx = -\frac{x}{4} \sqrt{X^3} + \frac{a^2}{8} \left( x \sqrt{X} + a^2 \arcsin \frac{x}{a} \right)\)\\[8pt]

            \(\int \sin(ax) \, dx = -\frac{1}{a} \cos(ax)\) \\[8pt]
            \(\int \sin^2(ax) \, dx = \frac{1}{2}x - \frac{1}{4a} \sin(2ax)\) \\[8pt]
            \(\int x \sin(ax) \, dx = \frac{\sin(ax)}{a^2} - \frac{x \cos(ax)}{a}\) \\[8pt]
            \(\int x^2 \sin(ax) \, dx = \frac{2x}{a^2} \sin(ax) - \left( \frac{x^2}{a} - \frac{2}{a^3} \right) \cos(ax)\) \\[8pt]

            \(\int \cos(ax) \, dx = \frac{1}{a} \sin(ax)\) \\[8pt]
            \(\int \cos^2(ax) \, dx = \frac{1}{2}x + \frac{1}{4a} \sin(2ax)\) \\[8pt]
            \(\int x \cos(ax) \, dx = \frac{\cos(ax)}{a^2} + \frac{x \sin(ax)}{a}\) \\[8pt]
            \(\int x^2 \cos(ax) \, dx = \frac{2x}{a^2} \cos(ax) + \left( \frac{x^2}{a} - \frac{2}{a^3} \right) \sin(ax)\) \\[8pt]

            \(\int x e^{a x} \, dx = \frac{e^{a x}}{a^2} \cdot (a x - 1)\)
      \end{minipage}

      \newpage

      \section*{Wahrscheinlichkeitsrechnung und Statistik}
         \noindent
         \begin{minipage}[t]{0.45\textwidth}
            \subsection*{Parameter einer Verteilung}
            \begin{itemize}
               \item \textbf{Erwartungswert (Mittelwert):} 
               \[
                  \mu = \mathbb{E}[X]
               \]
               \item \textbf{Varianz oder Streuung:} 
               \[
                  \sigma^2 = \text{Var}(X) = \mathbb{E}[(X - \mu)^2] = \mathbb{E}[X^2] - (\mathbb{E}[X])^2  % chktex 3
               \]
               \item \textbf{Standardabweichung:}
               \[
                  \sigma = \sqrt{\text{Var}(X)}
               \]
            \end{itemize}
         \end{minipage}
         \hfill
         \begin{minipage}[t]{0.45\textwidth}
            \subsection*{Rechenregeln für Erwartungswerte}
               \begin{itemize}

                  \item \textbf{Linearität des Erwartungswertes:}
                  \[
                  \mathbb{E}[aX + bY] = a \mathbb{E}[X] + b \mathbb{E}[Y]
                  \]

                  \item \textbf{Erwartungswert einer Konstanten:}
                  \[
                  \mathbb{E}[c] = c, \quad \text{für eine Konstante } c
                  \]



                  \item \textbf{Multiplikation unabhängiger Zufallsvariablen:}
                  \[
                  \mathbb{E}[XY] = \mathbb{E}[X] \cdot \mathbb{E}[Y]
                  \]


                  \item \textbf{Erwartungswert einer Summe:}
                  \[
                  \mathbb{E}\left[\sum_{i=1}^n X_i\right] = \sum_{i=1}^n \mathbb{E}[X_i]
                  \]


                  \item \textbf{Quadratische Form:}
                  \[
                  \mathbb{E}[X^2] = \text{Var}(X) + (\mathbb{E}[X])^2   % chktex 3
                  \]


               \end{itemize}
         \end{minipage}
         \medskip
         \hrule

         \begin{center}
            \subsection*{Verteilungen}
         \end{center}

         \begin{table}[h!]
            \centering
            \renewcommand{\arraystretch}{1.5}
               \begin{tabular}{|l|c|c|c|c|c|}   % chktex 44
                  \hline   % chktex 44
                  & \textbf{\(W(X)\)} & \textbf{\(f(x)\)} & \textbf{\(E[X]\)} & \textbf{\(\text{Var}(X)\)} & \textbf{Diagramm der Dichte} \\ \hline % chktex 44
                  Gleichverteilung &
                  \([a, b]\) &
                  \(\displaystyle f(x) = 
                  \begin{cases} 
                     \frac{1}{b-a}, & a \leq x \leq b \\ 
                     0, & \text{sonst}
                  \end{cases}\) &
                  \(\frac{a+b}{2}\) &
                  \(\frac{(b-a)^2}{12}\) &   % chktex 3
                  \begin{tikzpicture}[scale=0.6]
                     \draw[->] (0,0) -- (4.5,0) node[right] {\(x\)};
                     \draw[->] (0,0) -- (0,2.5) node[above] {\(f(x)\)};
                     \draw[thick] (1,0) -- (1,2) -- (3,2) -- (3,0);
                     \node[below] at (1,0) {\(a\)};
                     \node[below] at (3,0) {\(b\)};
                     \node[right] at (3,2) {\(\frac{1}{b-a}\)};
                  \end{tikzpicture} \\ \hline   % chktex 31, % chktex 44

                  Normalverteilung &
                  \(\mathbb{R}\) &
                  \(\displaystyle f(x) = \frac{1}{\sqrt{2\pi\sigma^2}} e^{-\frac{(x-\mu)^2}{2\sigma^2}}\) & % chktex 3
                  \(\mu\) &
                  \(\sigma^2\) &
                  \begin{tikzpicture}[scale=0.6]
                     \draw[->] (-3.5,0) -- (3.5,0) node[right] {\(x\)};
                     \draw[->] (0,0) -- (0,2.5) node[above] {\(f(x)\)};
                     \draw[thick, domain=-3:3, samples=100] 
                     plot (\x,{2*exp(-\x*\x/2)});
                     \node[below] at (0,0) {\(\mu\)};
                     \node[right] at (0.5,2) {\(\frac{1}{\sigma\sqrt{2\pi}}\)};
                  \end{tikzpicture} \\ \hline   % chktex 31, % chktex 44
               \end{tabular}
               \caption*{\footnotesize Normalverteilung \(=\) Standardnormalverteilung \(=\) Gauß-Verteilung}
               \caption*{\footnotesize Gleichv. \(=\) Rechteckv. \(=\) Uniformv. \(=\) Konstante Verteilung \(=\) Diskrete Gleichv. \(=\) Stetige Gleichv.}
         \end{table}

         \begin{table}[h!]
            \centering
            \footnotesize
            \renewcommand{\arraystretch}{2} 
            \setlength{\tabcolsep}{9pt} 
            \begin{tabular}{|c|c|c|c|c|}  % chktex 44
            \hline   % chktex 44
            \textbf{Eigenschaft} & Verteilung & Parameter & Dichtefunktion \(f_x(X)\) & Verteilungsfunktion \(F_x(X)\) \\ \hline  % chktex 44
            \textbf{Beschreibung} & Normal- oder Gaußverteilung & \(\mu = 0, \sigma = 1\) & 
            \(\displaystyle f(x) = \frac{1}{\sqrt{2\pi}} e^{-\frac{x^2}{2}}\) & 
            \(\displaystyle F(x) = \frac{1}{\sqrt{2\pi}} \int_{-\infty}^x e^{-\frac{t^2}{2}} dt\) \\ \hline % chktex 44
            \end{tabular}
            \caption*{\footnotesize Normal- und Gaußverteilung beschreibt dasselbe (Synonyme)}
         \end{table}

\newpage

%%% Wahrscheinlichkeit und Zufallsvariablen =======================================================
\section*{\underline{Prüfungsaufgabe 1: Teil 1}}

\pagecolor{Aufgabe1Teil1}

   \section*{Wahrscheinlichkeit}

   \noindent
   \begin{minipage}[t]{0.45\textwidth}
      \subsection*{Rechengesetze}

      \underline{Additionssatz} \\[5pt]
         \(P(A \cup B) = P(A) + P(B) - P(A \cap B)\) \\[5pt]
         \(P(A \cup B \cup C) = P(A) + P(B) + P(C) - P(A \cap B) - P(A \cap C) - P(B \cap C) + P(A \cap B \cap C)\) \\

      \underline{Komplementärereignis} \\[5pt]
         \(P(\overline{A}) = 1 - P(A)\) \\[5pt]
         \(P(\overline{A} \mid B) = 1 - P(A \mid B)\) \\

      \underline{Zerlegung eines Ereignisses} \\[5pt]
         \(P(A) = P(A \cap B) + P(A \cap \overline{B})\) \\

      \underline{Kommutativgesetz} \\[5pt]
         \(A \cup B = B \cup A\) \\[5pt]
         \(A \cap B = B \cap A\) \\

      \underline{Assoziativgesetz} \\[5pt]
         \(A \cup (B \cup C) = (A \cup B) \cup C\) \\[5pt]
         \(A \cap (B \cap C) = (A \cap B) \cap C\) \\

      \underline{Distributivgesetz} \\[5pt]
         \(A \cup (B \cap C) = (A \cup B) \cap (A \cup C)\) \\[5pt]
         \(A \cap (B \cup C) = (A \cap B) \cup (A \cap C)\) \\

      \underline{De Morgan'sche Gesetze} \\[5pt]
         \(\overline{A \cap B} = \overline{A} \cup \overline{B}\) \\[5pt]
         \(\overline{A \cup B} = \overline{A} \cap \overline{B}\) \\

      \underline{Unabhängige Ereignisse} \\[5pt]
         \text{stochastisch \textbf{unabhängig}} \\
            \(P(A \cap B) = P(A) \cdot P(B)\)  \\[5pt]
         \text{stochastisch abhängig} \\[5pt]
            \(P(A \cap B) \neq P(A) \cdot P(B)\) \\
   \end{minipage}
   \hfill
   \begin{minipage}[t]{0.45\textwidth}

      \section*{Bedingte Wahrscheinlichkeit} 
         \(P(A \mid B) = \frac{P(A \cap B)}{P(B)}, \quad \text{falls } P(B) > 0\)

      \section*{Formel der totalen Wahrscheinlichkeit}
         \(P(A) = \sum_{i=1}^N P(A \mid B_i) \cdot P(B_i)\)

      \section*{Bayes}
         \textbf{falls \(P(B)\) nicht bekannt} \\[10pt]
            \(P(A_k \mid B) = \frac{P(B \mid A_k) \cdot P(A_k)}{\sum_{i=1}^N P(B \mid A_i) P(A_i)}\) \\[10pt]
         \textbf{falls \(P(B)\) bekannt} \\[10pt]
            \(P(A_k \mid B) = \frac{P(B \mid A_k) \cdot P(A_k)}{P(B)}\)

      \section*{Weitere Formeln}

         \textbf{Multiplikationsregel mit bedingter \\ Wahrscheinlichkeit} \\[10pt]
            \(P(A_1 \cap A_2 \mid B) = P(A_1 \mid B) \cdot P(A_2 \mid B)\) \\

         \textbf{Relative Häufigkeit} \\[5pt]
            \(n(A)\): Anzahl der Häufigkeit / Eintreten \\[5pt]
            \(N\): Gesamtanzahl der Versuche / Beobachtungen \\[5pt]
            \(P(A) = \frac{n(A)}{N}\)

      \section*{Baumdiagramm}
      \begin{tikzpicture}[sibling distance=3cm]
         \node {P\((\cdot)\)}
         child {node {A} 
         child {node {D} 
         child {node {X} edge from parent node [left][yshift = 0.2cm] {\(\frac{4}{7}\)}}
         child {node {Y = \(\overline{X}\)} edge from parent node [right][yshift = 0.2cm] {\(\frac{3}{7}\)}}
         edge from parent node [left][yshift = 0.2cm] {\(\frac{1}{2}\)}}
         child {node {E = \(\overline{D}\)} edge from parent node [right][yshift = 0.2cm] {\(\frac{1}{2}\)}}
         edge from parent node [left][yshift = 0.2cm] {\(\frac{3}{10}\)}
         } 
         child {node {B = \(\overline{A}\)} edge from parent node [right][yshift = 0.2cm] {\(\frac{7}{10}\)}};
      \end{tikzpicture}

   \end{minipage}
   
   
\newpage

%%% Prüfungsaufgabe 1: Teil 2 =====================================================================
\section*{\underline{Prüfungsaufgabe 1: Teil 2}}

\pagecolor{Aufgabe1Teil2}

\noindent
\begin{minipage}[t]{0.45\textwidth}

\section*{Diskrete Wahrscheinlichkeit}

   \subsection*{Wahrscheinlichkeitsdichtefunktion}

   \begin{enumerate}
      \item \(\int_{-\infty}^{\infty} f_X(x) \, dx = 1\)
      \item Wahrscheinlichkeit ist nicht negativ!
      \item \(f_X(x)\) ist stückweise kontinuierlich
   \end{enumerate}

   \(f_X(x) = \frac{dF_X(x)}{dx}\)

   \subsection*{Erwartungswert}

      \(\mu = E[X]\)\\[10pt]
      \(E[X] = \int_{-\infty}^{\infty} x \, f_X(x) \, dx\)  \\[10pt]
      \(E[X^2] = \int_{-\infty}^{\infty} x^2 \, f_X(x) \, dx\) \\[5pt]

      \textbf{Bedingter Erwartungswert}

      \begin{align*}
         E[X \mid X > \mathbf{a}]    &= \frac{1}{P(X > \mathbf{a})} \int_{\mathbf{a}}^{\infty} x f_X(x) \, dx \\[5pt]
         &= \frac{1}{1 - F_X(\mathbf{a})} \int_{\mathbf{a}}^{\text{Letzte Grenze}} x f_X(x) \, dx
      \end{align*}

   \subsection*{Varianz}

      \(\sigma^2 = \text{Var}(X) = E[X^2] - (E[X])^2\)  % chktex 3

   \end{minipage}
   \hfill
   \begin{minipage}[t]{0.45\textwidth}

   \subsection*{Verteilungsfunktion}

      \(F_X(-\infty) = 0\) \\[10pt]
      \(0 \leq F_X(x) \leq 1\) \\[10pt]
      \(F_X(\infty) = 1\) \\[10pt]
      \(F_X(x_1) = \int_{-\infty}^{x} f_X(t) \, dt\) \\[10pt]
      \(F_X(x_2) = \int_{-\infty}^{x} f_X(t) \, dt + F_X(x_1)\) \\[10pt]
      \(F_X(x) = P(\{X \leq x\}) = \int_{-\infty}^{x} f_X(t) \, dt\) \\[10pt]
      \(P(X > x) = \int_{x}^{\text{Grenze Ende}} f_X(x) \, dx\)
      \begin{align*}
         P(\{x_1 < X < x_2\})    &= F_X(x_2) - F_X(x_1) \\[10pt]
         &= \int_{x_1}^{x_2} f_X(t) \, dt
      \end{align*}
      \[
      F_X(x) = 
      \begin{cases}
         0, & \text{für } x > \text{Grenze 1} \\[5pt]
         \text{\ldots}, & \text{Grenze 1} \le x \le \text{Grenze 2} \\[5pt]
         1, & \text{für } x > \text{Grenze 2}
      \end{cases}
      \] 

      \(F_{XY}(x, y) = F_X(x) \cdot F_Y(y)\)
   \end{minipage}
\newpage

%%% Prüfungsaufgabe 2 und 3 =======================================================================

\section*{\underline{Prüfungsaufgabe 2 und 3}}

\pagecolor{Aufgabe2}

\tableofcontents
\newpage

\noindent

\section{Kovarianz}

   \subsection*{Eigenschaften der Kovarianz}
      \begin{itemize}
         \item \textbf{Unkorrelierte Variablen:} Falls \( c_{X_1X_2}(\kappa) = 0 \), sind \( X_1 \) und \( X_2 \) unkorreliert. Bedeutet also, dass keine lineare Beziehung zwischen den Variablen besteht.
         \item \textbf{Orthogonalität:} Falls \( r_{X_1 X_2} = 0 \), sind die Variablen orthogonal.
         \item \textbf{Unabhängigkeit:} Wenn \( X_1 \) und \( X_2 \) unabhängig sind, gilt:
         \[
            f_{X_1 X_2}(x_1, x_2) = f_{X_1}(x_1) f_{X_2}(x_2) \quad \text{und} \quad c_{X_1 X_2} = 0.
         \]
      \end{itemize}

\newpage

\section{SOMF, WSS}
   \subsection{Second-Order Moment Function (SOMF)}
      Die \textbf{Second-Order Moment Function (SOMF)} misst die Ähnlichkeit eines Zufallsprozesses mit einer zeitverschobenen Version desselben Prozesses.

      \begin{itemize}
         \item SOMF ist definiert als: 
         \[
         r_{XX}(n+k, n) = E[X(n+k) \cdot X(n)]
         \]
         \item Für einen komplexen Zufallsprozess \( X(n) \) ist die SOMF definiert als:
         \[
         r_{XX}(n_1, n_2) = E[X(n_1) \cdot X^*(n_2)].
         \]
      \end{itemize}

   \subsection{Wide-Sense Stationarity (WSS)}
      Ein Zufallsprozess \( X(n) \) ist \textbf{Wide-Sense Stationary (WSS)}, wenn:
      \begin{enumerate}
         \item Der Erwartungswert \( E[X(n)] \) konstant ist.
         \item Die SOMF nur von der Zeitdifferenz \( \kappa \) abhängt:
         \[
         r_{XX}(n+\kappa, n) = r_{XX}(\kappa)
         \]
      \end{enumerate}

   \subsection{Cross-Second-Order Moment Function}
      Die \textbf{Cross-SOMF} beschreibt die Ähnlichkeit zwischen zwei Zufallsprozessen \( X_1(n) \) und \( X_2(n) \) \textbf{IM NICHT STATIONÄREN ZUSTAND}:
      \[
         r_{X_1 X_2}(n_1, n_2) = E[X_1(n_1) \cdot X_2(n_2)].
      \]

      \begin{itemize}
         \item Wenn \( X_1(n) \) und \( X_2(n) \) gemeinsam WSS sind, hängt die Cross-SOMF nur von \( \kappa\) ab:
         \[
            r_{X_1 X_2}(\kappa) = E[X_1(n + \kappa) \cdot X_2(n)], \quad \text{und somit ergibt sich}
         \]
         \begin{center}
            \(c_{X_1 X_2}= r_{X_1 X_2}(\kappa) = E[X_1(n + \kappa) \cdot X_2(n)]\)
         \end{center} 
      \end{itemize}

   \subsection{Central Second-Order Moment Function (Kovarianzfunktion)}
      Die zentrale SOMF, auch \textbf{Kovarianzfunktion} genannt, wird definiert als:
      \[
         c_{XX}(n+\kappa, n) = E[(X(n+\kappa) - E[X(n+\kappa)]) \cdot (X(n) - E[X(n)])], \quad \text{bzw.}
      \] 
      \begin{center}
         \(c_{XX}(n+\kappa, n) = E[X(n+\kappa) \cdot X(n)] - E[X(n+\kappa)] \cdot E[X(n)]\), \quad \text{falls \(E[X(n)] = 0\)} \\[10pt]
      \(c_{XX}(n+\kappa, n) = r_{XX} = E[X(n+\kappa) \cdot X(n)]\)
      \end{center}
      \medskip
      \begin{itemize}
         \item Die Kovarianzfunktion kann durch die SOMF ausgedrückt werden:
         \[
         c_{XX}(n+\kappa, n) = r_{XX}(n+\kappa, n) - E[X(n+\kappa)]E[X(n)].
         \] \\[10pt]
         \item Für zwei Prozesse \( X(n) \) und \( Y(n) \) verwenden wir die \textbf{Kreuzkovarianzfunktion}:
         \[
         c_{XY}(n+\kappa, n) = r_{XY}(n+\kappa, n) - E[X(n+\kappa)]E[Y(n)],
         \]
         \[
         c_{XY}(n+\kappa, n) = E[X(n+k)Y(n)] - E[X(n+\kappa)]E[Y(n)].
         \] \\[5pt]
   \end{itemize}

   Des Weiteren gilt noch: \(c_{XY}(k) = c_{YX}(-k)\)


\newpage

\section{Leistungsspektraldichte (PSD) und Spektrum}
\subsection{Leistungsspektraldichte}

   \subsubsection{Weißes Rauschen}

   Ein \textbf{weißes Rauschen} ist ein Zufallsprozess \(X(n)\), dessen \textbf{Leistungsdichtespektrum (Power Spectral Density, PSD)} konstant über alle Frequenzen ist:
   \[
   S_{XX}(e^{j\omega}) = \sigma_X^2
   \]

   \subsubsection*{Eigenschaften}
      \begin{enumerate}
         \item \textbf{Konstante PSD:} Das PSD bleibt unabhängig von der Frequenz konstant. Dies bedeutet, dass das weiße Rauschen eine gleichmäßige Energieverteilung über alle Frequenzen hat.
         \item \textbf{Zweite Momentfunktion (SOMF):} Die SOMF für einen weißen Prozess wird durch die inverse Fourier-Transformation der PSD erhalten:
         \[
            r_{XX}(\kappa) = \sigma_X^2 \delta(\kappa)
         \]
            wobei \(\delta(\kappa)\) die Kronecker-Delta-Funktion ist, definiert als:
         \[
         \delta(\kappa) = 
         \begin{cases} 
            1, & \text{wenn } \kappa = 0 \\
            0, & \text{wenn } \kappa \neq 0
         \end{cases}
         \]
            \item \textbf{Statistik:} Die Zufallsvariablen eines weißen Rauschens sind unkorreliert, haben eine konstante Varianz \(\sigma_X^2\) und einen Erwartungswert von 0.
         \item Unabhängigkeit und identische Verteilung sind klassische Eigenschaften von Rauschprozessen.
      \end{enumerate}

   \subsubsection{Kreuzleistungsspektraldichte CPSD}
      Die \textbf{Cross-Power Spectral Density (CPSD)} beschreibt die Frequenzabhängigkeit der statistischen Korrelation zwischen zwei zufälligen Prozessen \(X(n)\) und \(Y(n)\). Dies ist ein wesentlicher Begriff in der Frequenzbereichsanalyse, insbesondere bei der Untersuchung von Beziehungen zwischen zwei Prozessen.

   \subsubsection*{Eigenschaften}
      \begin{enumerate}
      \item \textbf{Symmetrie}:
      \[
         S_{XY}(e^{j\omega}) = S_{YX}(e^{j\omega})^*  % chktex 3
      \]
      Wenn \(X(n)\) und \(Y(n)\) reellwertig sind, dann gilt:
      \[
         S_{XY}(e^{j\omega}) = S_{YX}(e^{-j\omega})
      \]

      \item \textbf{Reale und imaginäre Anteile}:
      \begin{itemize}
         \item Der reale Anteil von \(S_{XY}(e^{j\omega})\) ist gerade.
         \item Der imaginäre Anteil von \(S_{XY}(e^{j\omega})\) ist ungerade.
      \end{itemize}

      \item \textbf{Orthogonalität}:
      \[
         \text{Wenn } X(n) \text{ und } Y(n) \text{ orthogonal sind, dann gilt: } S_{XY}(e^{j\omega}) = 0
      \]
      \end{enumerate}

   \newpage

   \subsection{Definition des Spektrums}

      Das Spektrum eines stationären Zufallsprozesses \(X(n)\) wird definiert als die Fourier-Transformation der Kovarianzfunktion:
      \[
         C_{XX}(e^{j\omega}) = \sum_{n=-\infty}^{\infty} c_{XX}(n) e^{-j\omega n}
      \]

   \subsubsection{Eigenschaften}
      \begin{enumerate}
      \item Die Spektralfunktion \(C_{XX}(e^{j\omega})\) existiert und ist beschränkt sowie gleichmäßig stetig, wenn:
      \[
         \sum_{n} |c_{XX}(n)| < \infty
      \]
         \item Die Inversion ist gegeben durch:
      \[
         c_{XX}(n) = \frac{1}{2\pi} \int_{-\pi}^{\pi} C_{XX}(e^{j\omega}) e^{j\omega n} d\omega
      \]
      \item \(C_{XX}(e^{j\omega})\) ist:
      \begin{itemize}
         \item Reellwertig
         \item \(2\pi\)-periodisch
         \item Nicht-negativ
      \end{itemize}
      \end{enumerate}

   \subsubsection{Kreuzspektrum}
      Für zwei gemeinsam stationäre Zufallsprozesse \(X_1(n)\) und \(X_2(n)\) ist das Kreuzspektrum definiert als:
      \[
         C_{X_1X_2}(e^{j\omega}) = \sum_{n=-\infty}^\infty c_{X_1X_2}(n) e^{-j\omega n}
      \]
      Mit der Eigenschaft:
      \[
         C_{X_1X_2}(e^{j\omega}) = C_{X_2X_1}(e^{j\omega})^*   % chktex 3
      \]
      Die Inversion lautet:
      \[
         c_{X_1X_2}(n) = \frac{1}{2\pi} \int_{-\pi}^{\pi} C_{X_1X_2}(e^{j\omega}) e^{j\omega n} d\omega
      \]

   \subsubsection{Reellwertige Prozesse}
      Wenn \(X_1(n)\) und \(X_2(n)\) reellwertig sind, gilt:
      \begin{itemize}
         \item \(C_{X_1X_2}(e^{j\omega}) = C_{X_2X_1}(e^{-j\omega})\)
         \item \(C_{X_1X_2}(e^{j\omega}) = C_{X_1X_2}(e^{-j\omega})^* = C_{X_2X_1}(e^{-j\omega}) = C_{X_2X_1}(e^{j\omega})^*\)  % chktex 3
      \end{itemize}

\newpage
\section{FIR und IIR Filter}
   \subsection{Einführung Digitale Filterung}
   \subsubsection{Allgemeines Konzept der Filterung}
      \begin{itemize}
         \item \textbf{Filterung} dient der Reduktion oder Verstärkung bestimmter Signalanteile.
         \item Digitale Filter arbeiten mit \textbf{abgetasteten, zeitdiskreten Signalen} und lassen sich in zwei Haupttypen unterteilen:
         \begin{enumerate}
            \item \textbf{Finite Impulse Response (FIR):} Haben eine endliche Impulsantwort.
            \item \textbf{Infinite Impulse Response (IIR):} Besitzen eine unendliche Impulsantwort.
         \end{enumerate}
         \item Der Frequenzbereich eines Filters wird durch die Übertragungsfunktion \(H(e^{j\omega})\) beschrieben:
         \[
            H(e^{j\omega}) = \frac{Y(e^{j\omega})}{X(e^{j\omega})}
         \]
         wobei \(X(e^{j\omega})\) und \(Y(e^{j\omega})\) die Fourier-Transformationen des Eingangs- und Ausgangssignals sind.
      \end{itemize}

   \smallskip

   \subsubsection{Lineare Differenzengleichung mit konstanten Koeffizienten (LCCDE)}
      \begin{itemize}
         \item Zeitinvariante, lineare diskrete Systeme werden durch eine LCCDE beschrieben:
         \[
            \sum_{k=0}^{M} a_k y(n-k) = \sum_{r=0}^{N-1} b_r x(n-r)
         \]
         \item \textbf{FIR-System:} Hat \(M = 0\):
         \[
            y(n) = \sum_{r=0}^{\text{Ordnung}} b_r x(n-r)
         \] \\[5pt]
         \[
            y(n) = \sum_{r=0}^{\text{Verzögerung} - 1} b_r x(n-r)
         \]
         \item \textbf{IIR-System:} Hat \(M > 0\):
         \[
            \sum_{k=0}^{M} a_k y(n-k) = \sum_{r=0}^{N-1} b_r x(n-r)
         \]
      \end{itemize}


   \newpage
   \subsection{Endliche Impulsantwort: FIR - Filter}  % chktex 8
      \subsection*{Definition}
         FIR-Filter sind lineare zeitinvariante Systeme mit einer endlichen Impulsantwort.\\ 
         \\
         \textbf{Differenzengleichung}:
         \[
            y(n) = \sum_{r=0}^{Information} b_r x(n-r), \quad \text{Information: Ordnung O = Verzögerung N - 1} % chktex 8
         \]
         Entspricht der Faltungsformel:
         \[
            y(n) = \sum_{k=-\infty}^{\infty} h(k)x(n-k) = \sum_{k=0}^{N-1} h(k)x(n-k)
         \]
         wobei \(h(k)\) die Impulsantwort des Filters darstellt.

   \subsection*{Vergleich IIR und FIR}
      \begin{enumerate}
         \item \textbf{Stabilität}
         \begin{itemize}
            \item IIR-Filter erfordern zusätzliche Stabilitätskriterien.
            \item FIR-Filter sind per Definition stabil.
         \end{itemize}
         \item \textbf{Lineare Phase}
         \begin{itemize}
            \item Es ist schwierig, eine lineare Phase mit einem IIR-Filter zu erreichen.
            \item FIR-Filter ermöglichen eine perfekte lineare Phase.
         \end{itemize}
         \item \textbf{Effizienz}
         \begin{itemize}
            \item IIR-Filter erreichen Spezifikationen für den Amplitudengang effizienter.
            \item IIR-Filter sind weniger rechenaufwendig als FIR-Filter.
         \end{itemize}
      \end{enumerate}

   \newpage
   \subsection{Unendliche Impulsantwort: IIR - Filter}   % chktex 8
      \subsection*{Definition}
         Ein IIR-Filter hat eine unendliche Impulsantwort.\\
         \\
         \textbf{Differenzengleichung:}
         \[
            \sum_{k=0}^M a_k y(n-k) = \sum_{r=0}^{N-1} b_r x(n-r)
         \]
         Die Übertragungsfunktion ist rational und wird dargestellt als:
         \[
            H(z) = \frac{\sum_{r=0}^{N-1} b_r z^{-r}}{1 + \sum_{k=1}^M a_k z^{-k}}
         \]

   \subsection*{Eigenschaften}
      \begin{enumerate}
         \item \textbf{Kausalität und Stabilität:}
         \begin{itemize}
            \item Für die Stabilität müssen die Wurzeln des Nennerpolynoms innerhalb des Einheitskreises im \(z\)-Bereich liegen (\(|z| < 1\)).
         \end{itemize}
         \item \textbf{Reelle und kausale Impulsantwort \(h(n)\):}
         \begin{itemize}
            \item Die Einheitspulsantwort \(h(n)\) ist reell, kausal und stabil.
         \end{itemize}
      \end{enumerate}

   \subsection*{Vergleich IIR und FIR}
      \begin{enumerate}
         \item \textbf{Stabilität:}
         \begin{itemize}
            \item IIR-Filter erfordern zusätzliche Stabilitätskriterien.
            \item FIR-Filter sind per Definition stabil.
         \end{itemize}
         \item \textbf{Lineare Phase:}
         \begin{itemize}
            \item Es ist schwierig, eine lineare Phase mit einem IIR-Filter zu erreichen.
            \item FIR-Filter ermöglichen eine perfekte lineare Phase.
         \end{itemize}
         \item \textbf{Effizienz:}
         \begin{itemize}
            \item IIR-Filter erreichen Spezifikationen für den Amplitudengang effizienter.
            \item IIR-Filter sind weniger rechenaufwendig als FIR-Filter.
         \end{itemize}
      \end{enumerate}

   \newpage
   \subsection{Lineare Filterung}

      Die lineare Filterung beschreibt, wie ein stationärer Eingangssignalprozess \(X(n)\) durch ein lineares zeitinvariantes System (LTI-System) verarbeitet wird, um einen Ausgangssignalprozess \(Y(n)\) zu erzeugen. Das System wird durch seine Einheitspulsantwort \(h(n)\) charakterisiert. Die Beziehung zwischen Eingangs- und Ausgangssignal ist durch die Faltungsoperation gegeben:
      \[
         Y(n) = \sum_{\rho} h(\rho) X(n-\rho)
      \]

      Ein LTI-System ist \textbf{stabil}, wenn die folgende Bedingung erfüllt ist:
      \[
         \sum_\rho |h(\rho)| < \infty
      \]

   \subsubsection{Kovarianzfunktion und Spektrum der Ausgabe}
      Die Kovarianzfunktion \(c_{YY}(\kappa)\) des Ausgangsprozesses \(Y(n)\) eines stabilen linearen Systems mit der Impulsantwort \(h(n)\) ist definiert als:
      \[
         c_{YY}(\kappa) = E[Y(n+\kappa)Y(n)]
      \] \\
      Ersetzt man \(Y(n)\) durch die gefilterte Version des Eingangssignals \(X(n)\):
      \[
         Y(n) = \sum_{\rho} h(\rho) X(n-\rho),
      \]
      führt dies zur Kovarianzfunktion:
      \[
         c_{YY}(\kappa) = \sum_{\rho}\sum_{\iota} h(\rho) h(\iota)^* c_{XX}(\kappa - \rho + \iota),   % chktex 3
      \]
      wobei \(c_{XX}(\kappa)\) die Kovarianzfunktion des Eingangssignals \(X(n)\) ist.

   \subsubsection*{Spektrum des Ausgangssignals}
      Durch Fourier-Transformation der Kovarianzfunktion ergibt sich die Spektraldichte des Ausgangssignals:
      \[
         C_{YY}(e^{j\omega}) = |H(e^{j\omega})|^2 C_{XX}(e^{j\omega}),
      \]
      wobei:
      \begin{itemize}
         \item \(H(e^{j\omega})\) die Übertragungsfunktion des Filters ist,
         \item \(C_{XX}(e^{j\omega})\) die Spektraldichte des Eingangssignals beschreibt.
      \end{itemize}

   \subsubsection{Kreuzkovarianzfunktion und Kreuzspektrum}
      Die Kreuz-Kovarianzfunktion des Ausgangssignals \(Y(n)\) und Eingangssignals \(X(n)\) eines stabilen LTI-Systems mit Impulsantwort \(h(n)\) ist definiert als:
      \[
         c_{YX}(\kappa) = E[Y(n+\kappa)X(n)]
      \]
      Durch Einsetzen von \(Y(n) = \sum_{\rho} h(\rho) X(n-\rho)\) ergibt sich:
      \[
         c_{YX}(\kappa) = \sum_\rho h(\rho) c_{XX}(\kappa-\rho)
      \]
      wobei \(c_{XX}(\kappa)\) die Autokovarianzfunktion des Eingangssignals \(X(n)\) ist.

\newpage

\section{Gefilterte Rauschprozesse: AR, MA und ARMA}
   \subsection{Lineare Filterung von Zufallsprozessen plus Rauschen}
      Lineare Filterung von stationären Zufallsprozessen \(X(n)\) in Kombination mit einem Rauschsignal \(V(n)\). Die Annahmen sind:
      \begin{enumerate}
         \item \(X(n)\) und \(V(n)\) sind stationär und unkorreliert
         \item \(h(n)\), \(X(n)\) und \(V(n)\) sind reellwertig.
      \end{enumerate}

      Der Ausgang \(Y(n)\) des Systems wird berechnet durch:
      \[
         Y(n) = \sum_{\rho=-\infty}^\infty h(\rho)X(n-\rho) + V(n)
      \]

   \subsection{Autoregressives (AR) Modell}
      Hauptidee des Filters besteht darin, das weiße Spektrum des Eingangssignals nach Wunsch zu formen.
      \subsubsection{Definition}
         Der AR-Prozess \(Y(n)\) der Ordnung \(p\) (AR(\(p\))) wird durch die \textbf{Differenzengleichung} definiert:  % chktex 36
         \[
            Y(n) + \sum_{k=1}^p a_k Y(n-k) = X(n),
         \]
         \textbf{Übertragungsfunktion:}
         \[
            H(e^{j\omega}) = \frac{1}{1 + \sum_{k=1}^p a_k e^{-j\omega k}}.
         \]

   \subsection{Filter mit gleitendem Mittelwert (MA)}
      Das \textbf{Moving-Average} Modell beschreibt einen Zufallsprozess Y(n), bei dem der aktuelle Wert als Linearkombination aktueller und vorheriger Werte eines weißen Rauschens X(n) dargestellt wird. % chktex 36
      \\
      \textbf{Differenzengleichung:}
      \[
         Y(n) = X(n) + \sum_{l=1}^{q} b_l X(n-l),
      \]
      \\
      \textbf{Übertragungsfunktion:}

      Diese Gleichung kann in die Übertragungsfunktion transformiert werden:
      \[
         H(e^{j\omega}) = 1 + \sum_{l=1}^q b_l e^{-j\omega l}.
      \]

   \subsection{Autoregressives Modell des gleitenden Durchschnitts (ARMA)}
      Der ARMA-Prozess \(Y(n)\) der Ordnung \(p\) und \(q\) (ARMA(\(p,q\))) wird durch die \textbf{Differenzengleichung} definiert:   % chktex 36
      \[
         Y(n) + \sum_{k=1}^p a_k Y(n-k) = X(n) + \sum_{l=1}^q b_l X(n-l),
      \]
      \textbf{Übertragungsfunktion:}
      \[
         H(e^{j\omega}) = \frac{1 + \sum_{l=1}^q b_l e^{-j\omega l}}{1 + \sum_{k=1}^p a_k e^{-j\omega k}}.
      \]

\newpage

   \section{Optimale lineare Systeme}

      \subsection{Wiener Filter}
         \subsubsection{Hauptkonzept}
            Der \textbf{Wiener Filter} ist ein lineares, zeitinvariantes Filter, das die Fehlerquadrate zwischen dem geschätzten Signal und dem tatsächlichen Signal minimiert. Dies basiert auf dem \textbf{Minimum Mean-Squared Error (MSE)-Kriterium}.

            \paragraph{Minimum Mean Squared Error (MSE)}
            \[
               q(h) = E\left[ \epsilon^2(n) \right].
            \]

            \[
               q(h) = E\left[\left(X(n) - \sum_{m} h(m)Y(n-m)\right)^2\right].   % chktex 3
            \]
            \\
            Der Fehler \( \epsilon(n) \) ist definiert als:
            \[
               \epsilon(n) = X(n) - \hat{X}(n), \quad \hat{X}(n) = \sum_{m} h_{opt}(m) Y(n-m).
            \]
            Hierbei ist:
            \begin{itemize}
               \item \( q(h)\): Der mittlere quadratische Fehler
               \item \( X(n) \): Das ursprüngliche Signal
               \item \( \hat{X}(n)\): Die geschätzte Version des ursprünglichen Signals
               \item \( Y(n) \): Das beobachtete Signal (Signal mit Rauschen überlagert)
               \item \( h(n) \): Die Impulsantwort des Filters
               \item \( \epsilon(n) \): Der Fehler
            \end{itemize}

         \subsubsection{Optimale Lösung \textbf{-} Optimierung}

            \textbf{Orthogonalitätsprinzip} \\
            Die Fehlerkomponente \(\epsilon(n)\) ist orthogonal zur beobachteten Komponente \(Y(n)\):
            \[
               c_{\epsilon Y}(\kappa) = E\left[ \epsilon(n + \kappa) Y(n) \right] = 0, \quad \forall \kappa \in \mathbb{Z}
            \]
            \\
            \textbf{Optimale Filterantwort} \\
            Im \textbf{Zeitbereich} ergibt sich die \textit{Wiener-Hopf-Gleichung}
            \[
               c_{XY}(\kappa) = \sum_{m =- \infty}^{\infty} h_{\text{opt}}(m) \ast c_{YY}(\kappa) 
            \]
            Im \textbf{Frequenzbereich} ergibt sich die optimale Filterantwort \( H_{\text{opt}}(e^{j\omega}) \)

            \[
               C_{XY}(e^{j\omega}) = H_{\text{opt}}(e^{j\omega}) \cdot C_{YY}(e^{j\omega}) 
            \]

            \[
               H_{\text{opt}}(e^{j\omega}) = \frac{C_{XY}(e^{j\omega})}{C_{YY}(e^{j\omega})},
            \]

\newpage

%%% Methode der kleinsten Quadrate ================================================================
\section*{\underline{Prüfungsaufgabe 4}}
\section*{Methode der kleinsten Quadrate}

\pagecolor{Aufgabe4}

\noindent
   \begin{minipage}[t]{0.45\textwidth}

      \subsection*{Grundlagen Matrizen}

         \[
         A = \begin{bmatrix}
            a_1 & a_2 \\
            b_1 & b_2
         \end{bmatrix}
         \]

         \[
            A^T = \begin{bmatrix}
               a_1 & b_1 \\
               a_2 & b_2
         \end{bmatrix}
         \]

         \[
            A^{-1} = \frac{1}{\det(A)} \begin{pmatrix} 
               d & -b \\ 
               -c & a 
         \end{pmatrix} = \frac{1}{ad - bc} 
         \begin{pmatrix} 
            d & -b \\ 
            -c & a 
         \end{pmatrix}
         \] \\

         \textbf{1.}
         \[
            X^T X = \begin{pmatrix}
               \sum_{i=1}^{N} x_i^2 & \sum_{i=1}^{N} x_i \\[15pt]
               \sum_{i=1}^{N} x_i & \sum_{i=1}^{N} 1 
         \end{pmatrix}
         \] \\

         \textbf{2.}
         \[
            X^T y = \begin{pmatrix}
               \sum_{i}^{N} x_i y_i \\[10pt]
               \sum_{i}^{N} y_i 
            \end{pmatrix}
         \] \\

         \textbf{3.}
         \[
            \hat{\theta} = (X^T X)^{-1} X^T y   % chktex 3
         \]

         \[
            \hat{\theta} = \begin{pmatrix}
               \hat{\theta} \\ 
               \hat{\theta_0}
         \end{pmatrix} \Rightarrow \begin{matrix}
            \hat{\theta} = x_2 \\ 
            \hat{\theta_0} = x_1
         \end{matrix} = \hat{\theta} = \begin{pmatrix}
            x_2 \\
            x_1
         \end{pmatrix}
         \] \\

         \textbf{4.}
         \[
            \hat{y} = X \hat{\theta} + \epsilon(x)= \begin{pmatrix}
               y_1 \\
               y_2 \\
               y_3 \\
               y_4 \\
               y_5 \\
               \ldots
            \end{pmatrix} = \begin{pmatrix}
            x_1 & 1 \\
            x_2 & 1 \\
            x_3 & 1 \\
            x_4 & 1 \\
            x_5 & 1 \\
            \ldots & \ldots
            \end{pmatrix}   \begin{pmatrix}
            \hat{\theta} \\
            \hat{\theta_0}
         \end{pmatrix} + \epsilon(x)
         \] 
   \end{minipage}
      \hfill
   \begin{minipage}[t]{0.45\textwidth}

         \subsection*{Weitere Formeln}

            \textbf{quadratische Abweichung} \\ 
            \textbf{(quadratischer Fehler)}

            \[
               \epsilon = \sum_{i = 1}^{N} (y_i - \hat{y}_i)^2 % chktex 3
            \] \\

            \textbf{mittlere quadratische Abweichung} \\ 
            \textbf{(mittlerer quadratischer Fehler)}
            \[
               \epsilon = \frac{1}{N} \sum_{i = 1}^{N} (y_i - \hat{y}_i)^2 % chktex 3
            \] \\

            \textbf{1. Stichprobenmittelwert für \(\bar{x}\) und \(\bar{y}\)} 
            \[
               \bar{x} = \frac{1}{N} \sum_{i=1}^{N} x_i, \quad \bar{y} = \frac{1}{N} \sum_{i=1}^{N} y_i
            \] \\

            \textbf{2. Steigung \(\hat{\theta}\)} \\
            \[
               \hat{\theta} = \frac{\sum_{i=1}^{N} (x_i - \bar{x})(y_i - \bar{y})}{\sum_{n = 1}^{N} (x_i - \bar{x})^2}  % chktex 3
            \] \\

            \textbf{3. Achsenabschnitt}
            \[
               \hat{\theta}_0 = \bar{y} - \hat{\theta} \bar{x}
            \] \\

            \textbf{4. Endformel}
            \[
               f(x) = \hat{y} = \hat{\theta}x_i + \hat{\theta}_0
            \] \\
   
   \end{minipage}

\end{document}